\chapter{Barriers and Digitals}\label{ch:dv:barriers-and-digitals}

A desk has sold a \omterm{def:dv:barriers-and-digitals:digital}{digital option} that pays 10 million if the index closes above 5\,000. One minute before the
close the index is at 4\,999.9. The option is worth 4.9 million, and a move of one index point in either
direction changes that by 1.25 million. To hedge it, the desk would need an index position equivalent to
6.25 billion, bought or sold in the last minute of trading and reversed if the index ticks back. No desk does
that. The digital's discontinuity makes it unhedgeable at the strike near expiry. The same discontinuity sits
at the level of every barrier option, which the market trades in large size in currencies and inside
structured products. This chapter prices digitals and barriers and explains how desks hedge what cannot be
delta-hedged, with an overhedge, with static portfolios of vanillas, and with a shifted barrier for discrete
monitoring. It ends with the risk that none of these remove, a price that jumps over the barrier.

\section{Digitals and their hedge}

\begin{definition}[Digital option]\label{def:dv:barriers-and-digitals:digital}
A \emph{digital option}\index{digital option} (binary option) pays a fixed amount if the underlying ends above
(call) or below (put) the strike, and nothing otherwise. Under Black--Scholes a cash digital call paying 1 is
worth $e^{-rT}\Phi(d_2)$.
\end{definition}

A digital is the limit of call spreads, $\bigl(C(K-\varepsilon)-C(K)\bigr)/\varepsilon$ as $\varepsilon\to0$,
so its price is minus the strike derivative of the call price. With a smile the call's volatility depends on the
strike, and differentiating gives
\[
D(K)=-\frac{\partial C}{\partial K}=D_{\mathrm{BS}}(K,\sigma(K))-\mathcal V(K)\,\frac{\partial\sigma}{\partial K}.
\]
With a downside skew $\partial\sigma/\partial K<0$, so the digital call is worth more than its Black--Scholes price
at the strike's volatility. The skew term is often the larger part of the difference between desks' quotes.

\begin{example}[A three-month digital on chapter 9's surface]\label{ex:dv:barriers-and-digitals:smile}
Scale chapter 9's surface to an index at 5\,000. A three-month at-the-money digital call on 10 million is worth
4.84 million with Black--Scholes at the at-the-money volatility of 16.4\%. With the smile it is worth 5.77 million:
the \omterm{def:dv:greeks-and-the-hedging-pnl:greeks}{vega} term adds 0.93 million, almost a fifth of the price.
\end{example}

Near expiry the digital's delta, $e^{-rT}\varphi(d_2)/(S\sigma\sqrt T)$ per unit paid, grows like $1/\sqrt T$ at the
strike and vanishes away from it (\cref{fig:dv:barriers-and-digitals:digital}, left). Its gamma changes sign at the
strike. A desk short a digital near the strike is either short or long a huge convexity, depending on which side
of the strike the index sits, and it cannot trade fast enough to follow. The practical answer is to replace the
digital by a call spread.

\begin{definition}[Call-spread overhedge]\label{def:dv:barriers-and-digitals:overhedge}
The \emph{call-spread overhedge}\index{call-spread overhedge} prices and hedges a sold digital call as the call
spread that dominates it: $N/\varepsilon$ calls struck at $K-\varepsilon$ less as many struck at $K$. The spread
pays at least the digital everywhere, and the difference in price, the overhedge, is charged to the buyer.
\end{definition}

The width $\varepsilon$ trades price against risk. A narrow spread is cheap but keeps most of the digital's
convexity: its delta can reach $N/\varepsilon$ per index point. A wide spread is easy to hedge but charges the client
for a payoff he did not buy. A desk sets the width from the largest loss it will accept per point of the
settlement price, the point at which the spread's slope becomes the risk. On the example, widths of 50, 100, 200
and 400 index points cost 0.25, 0.48, 0.89 and 1.56 million over the digital's 5.77 million.

\begin{omfigure}
\begin{tikzpicture}
\begin{groupplot}[group style={group size=2 by 1,horizontal sep=1.6cm},
  width=6.6cm,height=5cm,
  tick label style={font=\small},label style={font=\small},grid=major,grid style={gray!25},table/col sep=comma]
\nextgroupplot[xlabel={index},ylabel={delta (million per point)},xmin=4900,xmax=5100,ymin=0,ymax=0.18,
  xtick={4900,5000,5100},yticklabel style={/pgf/number format/fixed,/pgf/number format/precision=2},
  title={\small delta of a 10-million digital},
  legend columns=3,legend style={font=\footnotesize,at={(0.5,-0.32)},anchor=north,draw=none,fill=none}]
\addplot[omProp,thick,dashed] table[x=s,y=week]{figdata/derivatives/15-barriers-and-digitals/digital_delta.csv};
\addplot[omDef,thick] table[x=s,y=day]{figdata/derivatives/15-barriers-and-digitals/digital_delta.csv};
\addplot[omThm,very thick] table[x=s,y=hour]{figdata/derivatives/15-barriers-and-digitals/digital_delta.csv};
\legend{1 week,1 day,1 hour}
\nextgroupplot[xlabel={index at expiry},ylabel={payoff (million)},xmin=4600,xmax=5200,ymin=-0.5,ymax=11,
  xtick={4600,4800,5000,5200},
  title={\small digital and overhedge},
  legend columns=1,legend cell align=left,legend style={font=\footnotesize,at={(0.5,-0.32)},anchor=north,draw=none,fill=none}]
\addplot[omThm,very thick] table[x=s,y=digital]{figdata/derivatives/15-barriers-and-digitals/overhedge.csv};
\addplot[omDef,thick,dashed] table[x=s,y=spread]{figdata/derivatives/15-barriers-and-digitals/overhedge.csv};
\legend{digital,{call spread 4\,800--5\,000}}
\end{groupplot}
\end{tikzpicture}

\omcaption{Left: the delta of a 10-million digital struck at 5\,000 (20\% volatility) as expiry approaches; one
minute before the close it reaches 1.25 million per point, off this scale. Right: the digital's payoff and the
call spread that overhedges it; the extra payoff between 4\,800 and 5\,000 is what the overhedge charges for. Data: the chapter's
code.}\label{fig:dv:barriers-and-digitals:digital}
\end{omfigure}

\section{Barriers by reflection}

A barrier option (One Quant Book 2, chapter 19) switches on or off when the underlying touches a level $H$. With
continuous monitoring and Black--Scholes dynamics, all eight single barriers have closed forms, and all follow from
the reflection principle (One Quant Book 4, chapter 2).

\begin{proposition}[The down-and-out call by images]\label{prop:dv:barriers-and-digitals:image}
For $K\ge H$, carry $b=r-q$ and $\nu=b-\frac12\sigma^2$, the down-and-out call is
\[
C_{\mathrm{DO}}(S)=C_{\mathrm{BS}}(S,K)-\Bigl(\frac HS\Bigr)^{2\nu/\sigma^2}C_{\mathrm{BS}}\Bigl(\frac{H^2}S,K\Bigr).
\]
\end{proposition}
\begin{proof}
The function $u(S,t)=(H/S)^{2\nu/\sigma^2}C_{\mathrm{BS}}(H^2/S,K,t)$ solves the \omterm{def:dv:black-scholes-three-ways:pde}{Black--Scholes equation}, which can be
checked by substitution. It equals $C_{\mathrm{BS}}(H,K)$ on $S=H$, and its terminal payoff vanishes for $S>H$ when
$K\ge H$, because $H^2/S<H\le K$ there. The difference is therefore the solution that pays the call above the barrier
and zero on it. In the driftless log-price this is the method of images: the path density killed at $H$ is the free
density less its reflection in $H$.
\end{proof}

The other seven, with their in-out parity $V_{\mathrm{in}}+V_{\mathrm{out}}=V_{\mathrm{vanilla}}$ and rebates, are
combinations of four terms of this kind. The build implements all eight and tests parity for every case.

\begin{definition}[Barrier rebate]\label{def:dv:barriers-and-digitals:rebate}
A \emph{barrier rebate}\index{barrier rebate} is a fixed amount paid to the holder of a knock-out option when the
barrier is hit, or to the holder of a knock-in option at expiry if it never knocked in.
\end{definition}

\begin{definition}[One-touch option]\label{def:dv:barriers-and-digitals:onetouch}
A \emph{one-touch option}\index{one-touch option} pays a fixed amount if the underlying touches the barrier before
expiry, either at the hit or at expiry.
\end{definition}

\begin{definition}[No-touch option]\label{def:dv:barriers-and-digitals:notouch}
A \emph{no-touch option}\index{no-touch option} pays a fixed amount at expiry if the underlying never touches the
barrier. A one-touch paid at expiry and a no-touch on the same barrier add up to a zero-coupon bond.
\end{definition}

\begin{definition}[Double-barrier option]\label{def:dv:barriers-and-digitals:double}
A \emph{double-barrier option}\index{double-barrier option} has a lower and an upper barrier, and knocks in or out
when either is touched. The double no-touch, which pays if neither is touched, is its simplest form.
\end{definition}

Touch options are digitals on the path's extremes. The probability of touching a lower barrier $H<S$ by $T$ is
$\Phi\bigl((x-\nu T)/\sigma\sqrt T\bigr)+e^{2\nu x/\sigma^2}\Phi\bigl((x+\nu T)/\sigma\sqrt T\bigr)$, with
$x=\ln(H/S)$. A double barrier has no finite formula. Its killed density is a sine series in the log-price, whose
terms decay like $e^{-k^2\pi^2\sigma^2T/(2a^2)}$ for a corridor of log-width $a$, so a few dozen terms suffice.

\begin{example}[A barrier sheet]\label{ex:dv:barriers-and-digitals:sheet}
Spot and strike 100, one year, $r=3\%$, $q=1\%$, volatility 20\%. The call is worth 8.83. With a knock-out at 90 it is
worth 7.23, and the down-and-in 1.60, which adds up. A rebate of 2 paid at the hit raises the knock-out to 8.41, since
the one-touch at 90 paid at the hit is worth 0.592 (0.581 paid at expiry). The no-touch at 90 is worth 0.390, the up-and-out
call with barrier 130 is worth 3.10, and the double no-touch between 85 and 115 is worth 0.141
(\cref{fig:dv:barriers-and-digitals:barriers}, left).
\end{example}

\section{Static hedging}

A barrier's delta is discontinuous at the barrier, as the digital's is at its strike. The alternative to dynamic
hedging is a portfolio of vanillas that matches the barrier's value on the barrier and its payoff at expiry, so
that it can be unwound at no cost when the barrier is hit.

\begin{definition}[Static hedging]\label{def:dv:barriers-and-digitals:static}
\emph{Static hedging}\index{static hedging} replicates an exotic option with a portfolio of vanilla options fixed at
inception, traded again only at events such as a barrier hit or expiry, rather than rebalanced continuously.
\end{definition}

Carr, Ellis and Gupta's put--call symmetry gives the simplest static hedges. With zero carry and no smile, a call
struck at $K$ is worth, on the barrier $H$, exactly as much as $K/H$ puts struck at $H^2/K$, at every time. A
down-and-in call is therefore hedged by holding $K/H$ puts at $H^2/K$. If the barrier is hit, the desk swaps them for
the call at no cost. If it is not, the puts, whose strike is below the barrier, expire worthless, as the knock-in does.
With strike 100, barrier 90, one year and 20\% volatility, the hedge is 1.11 puts struck at 81, worth 1.498, the
closed-form price exactly.

The symmetry needs a symmetric smile, and an index smile is not symmetric. Suppose the barrier is hit with six months
left and chapter 9's six-month smile keeps its shape in moneyness. The 1.11 puts struck at 81 are then worth 2.08 and
the call struck at 100 only 0.83. The downside skew prices the puts, which sit below the barrier, at higher
volatilities than the call above it. A static hedge built from the symmetry would have cost far more than the option.
Skew breaks every model-free static hedge of this kind. Derman, Ergener and Kani's calendar hedge instead holds
options of several expiries chosen, under a model, so that the portfolio is worth zero on the barrier at a set of
dates.

\begin{example}[A calendar hedge of an up-and-out call]\label{ex:dv:barriers-and-digitals:calendar}
Strike 100, barrier 120, one year, 20\%, zero carry: the up-and-out call is worth 1.105. The terminal payoff is a call
at 100 less a call at 120 less 20 digitals at 120. Working back over $n$ equally spaced dates, the desk adds calls
struck at 120 of each expiry in the amounts that zero the portfolio on the barrier at each date. With 4 dates the hedge
costs 1.77 and is worth up to 5.4 on the barrier between dates, over the first nine months. With 16 dates it costs 1.28
and is worth at most 0.16. With 64 it costs 1.15 (\cref{fig:dv:barriers-and-digitals:static}). The error halves each time
the dates double. The last interval before expiry is the hardest, because there the payoff on the barrier jumps from
20 to 0.
\end{example}

\begin{omfigure}
\begin{tikzpicture}
\begin{groupplot}[group style={group size=2 by 1,horizontal sep=1.6cm},
  width=6.6cm,height=5cm,
  tick label style={font=\small},label style={font=\small},grid=major,grid style={gray!25},table/col sep=comma]
\nextgroupplot[xlabel={time (years)},ylabel={hedge value on the barrier},xmin=0,xmax=1,ymin=-2,ymax=12,
  title={\small calendar hedge on the barrier},
  legend columns=2,legend style={font=\footnotesize,at={(0.5,-0.32)},anchor=north,draw=none,fill=none}]
\addplot[omDef,thick,dashed] table[x=t,y=n4]{figdata/derivatives/15-barriers-and-digitals/calendar_barrier.csv};
\addplot[omThm,very thick] table[x=t,y=n16]{figdata/derivatives/15-barriers-and-digitals/calendar_barrier.csv};
\legend{4 dates,16 dates}
\nextgroupplot[xmode=log,log basis x=2,xlabel={hedge dates $n$},ylabel={hedge cost},xmin=1.5,xmax=90,ymin=1,ymax=2.5,
  xtick={2,4,8,16,32,64},xticklabels={2,4,8,16,32,64},
  title={\small cost against the closed form},
  legend columns=1,legend cell align=left,legend style={font=\footnotesize,at={(0.5,-0.32)},anchor=north,draw=none,fill=none}]
\addplot[omThm,very thick,mark=*,mark size=1.6pt] table[x=n,y=price]{figdata/derivatives/15-barriers-and-digitals/calendar_prices.csv};
\addplot[gray,thick,dashed] table[x=n,y=uoc]{figdata/derivatives/15-barriers-and-digitals/calendar_prices.csv};
\legend{static hedge,up-and-out call (1.105)}
\end{groupplot}
\end{tikzpicture}

\omcaption{A static hedge of an up-and-out call (strike 100, barrier 120, one year, 20\%) from calls of several
expiries. Left: the hedge's value on the barrier through time: zero at the dates it is built on, and large in the last
interval, where the payoff jumps. Right: its cost converges to the option's price as the dates multiply. Data: the
tutorial.}\label{fig:dv:barriers-and-digitals:static}
\end{omfigure}

\section{Discrete monitoring and the barrier shift}

Most barrier contracts outside currencies are monitored at discrete dates: daily closes or fixings. A discretely
monitored knock-out survives paths that dipped through the barrier between observations, so it is worth more than
the continuous formula says. Broadie, Glasserman and Kou showed that the continuous formula becomes accurate again if
the barrier is moved away from the spot.

\begin{definition}[Barrier shift]\label{def:dv:barriers-and-digitals:shift}
The \emph{barrier shift}\index{barrier shift} prices an option monitored every $\Delta t$ with the continuous
formula at the barrier $H\,e^{\pm\beta\sigma\sqrt{\Delta t}}$, moved away from the spot, with $\beta\approx0.5826$.
\end{definition}

The constant is $-\zeta(\frac12)/\sqrt{2\pi}$, the expected overshoot of a Gaussian random walk over a level, in units
of its step. The build keeps it in one named constant.

\begin{example}[How good is the shift]\label{ex:dv:barriers-and-digitals:shift}
The knock-out call of \cref{ex:dv:barriers-and-digitals:sheet} (barrier 90) is worth 7.23 with continuous monitoring. With
daily monitoring, a Monte Carlo of 200\,000 paths gives 7.48 (standard error 0.03), and the shifted formula 7.45. With
weekly monitoring they give 7.62 and 7.68, with monthly 8.01 and 8.04, and with four dates 8.33 and 8.38
(\cref{fig:dv:barriers-and-digitals:barriers}, right). For an up-and-out call with barrier 120, where the payoff is large
at the barrier, the shift is good daily (1.26 against 1.28) and poor quarterly (2.19 against 2.50). The correction is an
expansion in $\sqrt{\Delta t}$ and fails when the barrier is within a few steps of the spot or the payoff at the barrier is
large.
\end{example}

\begin{omfigure}
\begin{tikzpicture}
\begin{groupplot}[group style={group size=2 by 1,horizontal sep=1.6cm},
  width=6.6cm,height=5cm,
  tick label style={font=\small},label style={font=\small},grid=major,grid style={gray!25},table/col sep=comma]
\nextgroupplot[xlabel={spot},ylabel={price},xmin=88,xmax=130,ymin=0,ymax=32,
  title={\small barrier 90, strike 100},
  legend columns=3,legend style={font=\footnotesize,at={(0.5,-0.32)},anchor=north,draw=none,fill=none}]
\addplot[gray,thick,dashed] table[x=s,y=call]{figdata/derivatives/15-barriers-and-digitals/barrier_curves.csv};
\addplot[omThm,very thick] table[x=s,y=doc]{figdata/derivatives/15-barriers-and-digitals/barrier_curves.csv};
\addplot[omDef,thick] table[x=s,y=dic]{figdata/derivatives/15-barriers-and-digitals/barrier_curves.csv};
\legend{call,down-and-out,down-and-in}
\nextgroupplot[xmode=log,xlabel={monitoring dates a year},ylabel={down-and-out call},xmin=3,xmax=320,ymin=7,ymax=8.6,
  xtick={4,12,52,252},xticklabels={4,12,52,252},
  title={\small discrete monitoring},
  legend columns=2,legend style={font=\footnotesize,at={(0.5,-0.32)},anchor=north,draw=none,fill=none}]
\addplot[black,only marks,mark=*,mark size=1.8pt] table[x=n,y=mc]{figdata/derivatives/15-barriers-and-digitals/monitoring.csv};
\addplot[omThm,very thick] table[x=n,y=shifted]{figdata/derivatives/15-barriers-and-digitals/monitoring.csv};
\addplot[gray,thick,dashed] table[x=n,y=cont]{figdata/derivatives/15-barriers-and-digitals/monitoring.csv};
\legend{Monte Carlo,shifted barrier,continuous}
\end{groupplot}
\end{tikzpicture}

\omcaption{Left: a call, its down-and-out and its down-and-in by spot (one year, 20\%); they add up. Right: the
down-and-out call monitored 4, 12, 52 and 252 times a year: Monte Carlo against the continuous formula and the formula
at the shifted barrier. Data: the tutorial.}\label{fig:dv:barriers-and-digitals:barriers}
\end{omfigure}

The flat-volatility formula has a larger error than monitoring when the smile is steep. Under chapter 9's \omterm{def:dv:local-volatility:model}{local
volatility}, the same knock-out (zero rates, one year, daily monitoring) is worth 6.15 by simulation, against 6.51 from
the shifted \omterm{thm:dv:black-scholes-three-ways:formula}{Black--Scholes formula} at the at-the-money volatility of 19.2\%. \omterm{def:dv:local-volatility:model}{Local volatility} is higher at low spots, so
the barrier is hit more often. How far the smile moves such prices depends on the smile's dynamics as well as its
level, \omterm{def:dv:local-volatility:model}{local volatility} against stochastic volatility (chapters 9 and 10). Chapter 20 prices currency barriers with the
market's own correction, built from the option's \omterm{def:dv:greeks-and-the-hedging-pnl:greeks}{vanna} and \omterm{def:dv:greeks-and-the-hedging-pnl:greeks}{volga}.

\section{Gap risk at the barrier}

Every hedge in this chapter assumes the price passes through the barrier: the static hedge is unwound at the barrier,
the dynamic hedge trades near it, the stop-loss order is filled at it. A price that jumps over the barrier defeats all
three.

\begin{definition}[Gap risk]\label{def:dv:barriers-and-digitals:gap}
\emph{Gap risk}\index{gap risk} is the risk that a price moves discontinuously across a level at which a position's value
or hedge changes, a barrier, a strike near expiry, or a stop, so that the hedge planned at that level is executed at a
worse price or not at all.
\end{definition}

The reverse knock-out, whose payoff is large at the barrier, carries it in the purest form.

\begin{example}[A gap through a reverse barrier]\label{ex:dv:barriers-and-digitals:gap}
An up-and-out call, strike 100, barrier 120, has one month left, and the spot is 119. It is worth 1.38 and its delta is
$-1.37$: a rise brings the barrier closer and destroys the payoff
(\cref{fig:dv:barriers-and-digitals:gap}). The seller hedges by selling 1.37 shares. If the spot moves to 119.5, the
hedged position's P\&L is $+0.005$. If it gaps to 125, the option dies, and the seller gains its 1.38, but the short
shares lose $1.37\times6=8.19$: a net loss of 6.81, five times the premium.
\end{example}

\begin{omfigure}
\begin{tikzpicture}
\begin{groupplot}[group style={group size=2 by 1,horizontal sep=1.6cm},
  width=6.6cm,height=5cm,
  tick label style={font=\small},label style={font=\small},grid=major,grid style={gray!25},table/col sep=comma]
\nextgroupplot[xlabel={spot},ylabel={value},xmin=95,xmax=121,ymin=0,ymax=9,
  title={\small up-and-out call, one month}]
\addplot[omThm,very thick] table[x=s,y=value]{figdata/derivatives/15-barriers-and-digitals/reverse.csv};
\draw[gray,dashed] (axis cs:120,0) -- (axis cs:120,9);
\nextgroupplot[xlabel={spot},ylabel={delta},xmin=95,xmax=121,ymin=-1.6,ymax=0.8,
  title={\small its delta}]
\addplot[omDef,very thick] table[x=s,y=delta]{figdata/derivatives/15-barriers-and-digitals/reverse.csv};
\draw[gray,dashed] (axis cs:120,-1.6) -- (axis cs:120,0.8);
\draw[gray] (axis cs:95,0) -- (axis cs:121,0);
\end{groupplot}
\end{tikzpicture}

\omcaption{A reverse knock-out near its barrier: strike 100, barrier 120 (dashed), one month left, 20\%. Its value
peaks well before the barrier and its delta turns strongly negative; a hedger short the option is short the underlying
there, and loses on any gap through the barrier. Data: the chapter's code.}\label{fig:dv:barriers-and-digitals:gap}
\end{omfigure}

Gaps are not only a matter of weekends and earnings. On 15 January 2015 the Swiss National Bank discontinued its minimum
exchange rate of 1.20 Swiss francs per euro. The ECB's euro reference rate for the franc was 1.2010 on the 14th and
1.0280 on the 15th, a fall of 14.4\% between two fixings, and 1.0128 on the 16th. A barrier anywhere in that range was
crossed in a single move, and a hedge planned at the barrier could only be executed wherever the market next traded. A
floor that a central bank defends looks like the safest barrier in the market, and when it fails, it fails by a gap. Desks price the risk with jump models
(chapter 13) and hold reserves for it (chapter 27).

\section{Tutorial: barriers and digitals}\label{tut:dv:barriers-and-digitals:tut}

\begin{tutorial}
\textbf{Goal.} Price a digital with and without a smile and its \omterm{def:dv:barriers-and-digitals:overhedge}{call-spread overhedge}; price barriers in closed form,
under discrete monitoring and under \omterm{def:dv:local-volatility:model}{local volatility}; build two static hedges; measure the gap. \textbf{End state:} the
four figures and the numbers of the weekend problem.
\begin{tutsteps}
\item \textbf{A digital under a smile} is minus the strike derivative of the call:
\omcode{code/firm/barrier/firm_barrier.py}{34}{38}{The smile-consistent digital.}\label{lst:dv:barriers-and-digitals:digital}
\item \textbf{The \omterm{def:dv:barriers-and-digitals:shift}{barrier shift} and its check.} The shifted barrier, and a Monte Carlo of discretely monitored barriers
with exact lognormal steps:
\omcode{code/firm/barrier/firm_barrier.py}{129}{150}{Barrier shift and discrete-monitoring Monte Carlo.}\label{lst:dv:barriers-and-digitals:shift}
\item \textbf{A double no-touch} by its sine series:
\omcode{code/firm/barrier/firm_barrier.py}{110}{125}{Double no-touch.}\label{lst:dv:barriers-and-digitals:dnt}
\item \textbf{Run} \texttt{dv\_barrier.digital\_prices()}, \texttt{overhedge(200)}, \texttt{monitoring()},
\texttt{local\_vol\_doc()}, \texttt{calendar\_example()}, \texttt{reverse\_barrier()} and \texttt{fig\_barrier.py}.
\end{tutsteps}
\textbf{What to change next.} Price the digital as a call spread centred on the strike and compare its risk with the
overhedge; replace the calls at the barrier in the calendar hedge by puts; move the up-and-out barrier to 110 and re-run
the gap.
\end{tutorial}

\section{Build: barriers and digitals}\label{bld:dv:barriers-and-digitals:barrier}

\begin{build}
\bfield{Purpose} The miniature firm's barrier and digital pricer: closed forms with rebates, touches, double no-touch,
the discrete-monitoring shift, and a static-hedge builder; the engine that chapter 18's autocallables and chapter 20's
currency barriers call.
\bfield{Interface} \texttt{digital}, \texttt{digital\_delta}, \texttt{digital\_smile(s, k, t, r, q, vol\_of\_k)};
\texttt{call\_spread(\ldots, width, notional, below)}; \texttt{barrier(s, k, h, t, r, q, vol, kind, right, rebate)} with
\texttt{kind} in down-out, down-in, up-out, up-in; \texttt{hit\_probability}, \texttt{one\_touch(\ldots, at\_hit)},
\texttt{no\_touch}, \texttt{double\_no\_touch}; \texttt{bgk\_shift}, \texttt{BGK\_BETA}; \texttt{mc\_barrier};
\texttt{symmetry\_down\_in\_call}; \texttt{calendar\_hedge}, \texttt{calendar\_hedge\_on\_barrier}.
\bfield{Rules} In and out always add up to the vanilla (tested for all eight cases); discrete monitoring is priced with
the shift or by simulation, never with the continuous formula; a digital is always priced with the smile's slope; a
digital sold is booked at its overhedge.
\bfield{Acceptance tests} \texttt{code/firm/barrier/tests/}: in-out parity for every kind, right and strike; a distant
barrier gives the vanilla; the shifted formula matches a daily Monte Carlo; touch identities, and the double no-touch
against its single-barrier limit and a simulation; the smile digital reduces to the flat one without skew, and call
spreads bracket the digital; the symmetry hedge prices the down-and-in exactly, and the calendar hedge converges.
\bfield{Stretch} Barriers under local and stochastic volatility by finite differences (chapter 22); window barriers;
the vanna--volga correction of chapter 20.
\end{build}

\begin{omsources}
\item M.~Broadie, P.~Glasserman and S.~Kou, ``A continuity correction for discrete barrier options'', \emph{Mathematical
Finance} 7(4) (1997) 325--349.
\item E.~Derman, D.~Ergener and I.~Kani, ``Static options replication'', \emph{Journal of Derivatives} 2(4) (1995) 78--95.
\item P.~Carr, K.~Ellis and V.~Gupta, ``Static hedging of exotic options'', \emph{Journal of Finance} 53(3) (1998)
1165--1190.
\item Swiss National Bank, press release, 15 January 2015; European Central Bank, euro foreign exchange reference rates.
\end{omsources}

\section{Exercises}

\begin{exercise}[$\star$]\label{exo:dv:barriers-and-digitals:1}
Show that a cash digital call and a cash digital put on the same strike and expiry add up to a zero-coupon bond. What
does that imply for their deltas?
\end{exercise}

\begin{exercise}[$\star$]\label{exo:dv:barriers-and-digitals:2}
Why does a downside skew make a digital call more expensive than its Black--Scholes price at the strike's volatility?
\end{exercise}

\begin{exercise}[$\star$]\label{exo:dv:barriers-and-digitals:3}
A knock-out call and its knock-in are priced at 7.23 and 1.70 while the vanilla is 8.83. What is wrong?
\end{exercise}

\begin{exercise}[$\star\star$]\label{exo:dv:barriers-and-digitals:4}
With zero carry, check the symmetry hedge numerically: price 1.11 puts struck at 81 and the down-and-in call with
barrier 90 at 20\% volatility. Why do the puts expire worthless whenever the barrier is not hit?
\end{exercise}

\begin{exercise}[$\star\star$]\label{exo:dv:barriers-and-digitals:5}
A barrier is monitored daily at 20\% volatility. By how much does the shift move a barrier at 90? At 120?
\end{exercise}

\begin{exercise}[$\star\star$]\label{exo:dv:barriers-and-digitals:6}
Explain why the up-and-out call of \cref{ex:dv:barriers-and-digitals:gap} has a negative delta near the barrier, and why
its value peaks below the barrier.
\end{exercise}

\begin{exercise}[$\star\star\star$]\label{exo:dv:barriers-and-digitals:7}
\emph{Coding.} Re-run the discrete-monitoring comparison for a down-and-out put with barrier 90. Does the shift work as
well as for the call?
\end{exercise}

\begin{exercise}[$\star\star\star$]\label{exo:dv:barriers-and-digitals:8}
\emph{Find the flaw.} ``Our barrier book is safe: every barrier has a stop-loss order resting at its level, so the
hedge is executed exactly when the option knocks.''
\end{exercise}

\section{Problem: The Digital at Expiry}

\begin{problem}[{Weekend problem --- how wide a call spread}]\label{pb:dv:barriers-and-digitals:1}
A client buys a three-month digital that pays 10 million if the index, now at 5\,000, closes above 5\,000 at expiry. The
desk prices on chapter 9's surface scaled to 5\,000 and will book the digital at a \omterm{def:dv:barriers-and-digitals:overhedge}{call-spread overhedge}. Its limit: a
one-point error in the settlement price may not move the hedge by more than 50\,000.

\medskip\noindent\textbf{Part I --- The digital.}
\begin{enumerate}
\item Give the digital's price with the at-the-money volatility and with the smile.
\item What is the \omterm{def:dv:greeks-and-the-hedging-pnl:greeks}{vega} term, and why is it positive?
\item One minute before expiry with the index at 4\,999.9 and 20\% volatility, give the digital's delta per point and as
index notional.
\item Why can that delta not be traded?
\item What is the digital's gamma just below and just above the strike near expiry?
\end{enumerate}

\medskip\noindent\textbf{Part II --- The overhedge.}
\begin{enumerate}[resume]
\item Which call spread dominates the digital, and why is it the one the seller buys?
\item What width meets the desk's limit?
\item Give the spread's price and the overhedge charged to the client.
\item Give the overhedge for widths of 50, 100 and 400 points.
\item What does the client receive, relative to the digital, if the index settles at 4\,900?
\end{enumerate}

\medskip\noindent\textbf{Part III --- Barriers.}
\begin{enumerate}[resume]
\item Give the one-year down-and-out call (strike 100, barrier 90, 20\%, $r=3\%$, $q=1\%$) and its knock-in.
\item Give it with daily monitoring by simulation and by the shift.
\item Give it under chapter 9's \omterm{def:dv:local-volatility:model}{local volatility} and explain the difference from the flat price.
\item Give the double no-touch between 85 and 115.
\item Price the one-touch at 90 paid at the hit and at expiry, and explain the difference.
\end{enumerate}

\medskip\noindent\textbf{Part IV --- Judgement.}
\begin{enumerate}[resume]
\item Why is the overhedge fair to the client?
\item When would you use a centred call spread instead?
\item How would you hedge a reverse knock-out near its barrier?
\item State the \emph{named result}: the call-spread width that caps the hedger's loss at 50\,000 per settlement point on
the 10-million digital, and the overhedge it charges the client.
\item In one sentence: what do digitals and barriers share?
\end{enumerate}
\end{problem}

\section{Interview questions}

\begin{interviewq}[$\star$ \iqroles{trader}]\label{iq:dv:barriers-and-digitals:1}
How do you price and hedge a \omterm{def:dv:barriers-and-digitals:digital}{digital option}?
\end{interviewq}

\begin{interviewq}[$\star$ \iqroles{researcher}]\label{iq:dv:barriers-and-digitals:2}
Derive the price of a down-and-out call with the reflection principle.
\end{interviewq}

\begin{interviewq}[$\star\star$ \iqroles{trader, researcher}]\label{iq:dv:barriers-and-digitals:3}
Why does the skew matter for a digital, and in which direction?
\end{interviewq}

\begin{interviewq}[$\star\star$ \iqroles{developer}]\label{iq:dv:barriers-and-digitals:4}
Your Monte Carlo prices a daily-monitored knock-out above the closed form. Is that a bug?
\end{interviewq}

\begin{interviewq}[$\star\star$ \iqroles{risk}]\label{iq:dv:barriers-and-digitals:5}
What is \omterm{def:dv:barriers-and-digitals:gap}{gap risk} on a barrier book, and how would you measure and limit it?
\end{interviewq}

\begin{interviewq}[$\star\star\star$ \iqroles{trader, researcher}]\label{iq:dv:barriers-and-digitals:6}
Construct a static hedge for a down-and-in call. When does it fail?
\end{interviewq}
